


%%%%%%%%%%%%%%%%%%%%%%%%%%%%%%%%%%%%%%%%%%%%%%%%%%%%%%%%%%%%%%%%%%%%%%%%%%%%%%%%%%%%%%%%%%%%
%%%%%%%%%%%%%%%%%%%%%%%%%%%   NOUVELLES COMMANDES : RACCOURCIS   %%%%%%%%%%%%%%%%%%%%%%%%%%%
%%%%%%%%%%%%%%%%%%%%%%%%%%%%%%%%%%%%%%%%%%%%%%%%%%%%%%%%%%%%%%%%%%%%%%%%%%%%%%%%%%%%%%%%%%%%
%-------------------------------------------------
%---- Ensemles : entiers, reels, complexes... ----
%-------------------------------------------------
\newcommand{\N}{\mathbb{N}}
\newcommand{\Z}{\mathbb{Z}}
\newcommand{\Q}{\mathbb{Q}}
\newcommand{\R}{\mathbb{R}}
\newcommand{\F}{\mathbb{F}}
\newcommand{\Rb}{{\rm I\!R}} % permet d'avoir R en gras
%\renewcommand{\C}{\mathbb{C}} %/!\ la command "\C" semble déjà définie dansn le package hyperref
%\newcommand{\C}{\mathbb{C}} %/!\ la command "\C" semble déjà définie dansn le package hyperref % -> pour compiler avec dvi to pdf !!



%------------------------------
%---- Raccourcis : espaces ----
%------------------------------
\newcommand{\eh}[1]{\hspace*{#1mm}}
\newcommand{\ehh}{\hspace*{2mm}}
\newcommand{\ehhh}{\hspace*{3mm}}
\newcommand{\ev}{\vspace*{1mm}}
\newcommand{\evv}{\vspace*{2mm}}
\newcommand{\evvv}{\vspace*{3mm}}
\newcommand{\para}{\hspace*{8mm}}
\newcommand{\eexo}{\vspace{8mm}}
\newcommand{\eexocourt}{\vspace{7mm}}
\newcommand{\eu}[1][2]{\hspace*{#1pt}}



%---------------------------------------------
%---- Raccourcis : symboles mathématiques ----
%---------------------------------------------
\newcommand{\ssi}[1][0]{\hspace{#1mm}\Leftrightarrow\hspace{#1mm}}
\newcommand{\implique}[1][0]{\hspace{#1mm}\Rightarrow\hspace{#1mm}}
\newcommand{\ega}[1][0]{\hspace{#1mm}=\hspace{#1mm}}
\newcommand{\norme}[1]{\parallel #1 \parallel}
\renewcommand{\degre}{\ensuremath{^\circ} }
\newcommand{\imply}{\textrm{\hspace{2mm}\LARGE{\raisebox{-0.5mm}{$\Rightarrow$}}\hspace{2mm}}}
\newcommand{\notimply}{\textrm{\hspace{2mm}\LARGE{\raisebox{-0.5mm}{$\psCancel[linecolor=red,linewidth=2pt,cancelType=x]{\Rightarrow}$}}\hspace{2mm}}}
\newcommand{\rec}{^r\!}
\newcommand{\pgdc}{\textrm{pgdc}}
\newcommand{\ppmc}{\textrm{ppmc}}
\newcommand{\lnit}{\textrm{\it ln}}
\newcommand{\fois}{\hspace*{-0.5mm}\cdot\hspace*{-0.5mm}}



%-------------------------------------
%---- Raccourcis : mise en formes ----
%-------------------------------------
\newcommand{\guillemets}[1]{\textsl{``\hspace{0.5pt}#1''\hspace{1pt}}}
\newcommand{\textbs}[1]{\textbf{\textsl{#1}}}
\newcommand{\textbsc}[1]{\textbf{\textsc{#1}}}
\newcommand{\pvirg}[1]{\hspace*{#1 mm};\hspace{#1mm}}
\newcommand{\virg}[1][1]{\hspace*{#1 mm},\hspace{#1mm}}
\newcommand{\et}[1]{\hspace{#1mm}\textrm{et}\hspace{#1mm}}
\newcommand{\ou}[1]{\hspace{#1mm}\textrm{ou}\hspace{#1mm}}
\newcommand{\compy}[1]{\textsl{\guillemets{#1}}}
\newcommand{\egah}[1]{\overset{\textcolor{red!40!gray}{\scriptscriptstyle #1}}{\hspace{2mm}=\hspace{2mm}}}
\newcommand{\egab}[1]{\underset{\textcolor{red!40!gray}{\scriptscriptstyle #1}}{\hspace{2mm}=\hspace{2mm}}}
\newcommand{\egahb}[2]{%
  \mathrel{\stackunder[2pt]{\stackon[4pt]{\hspace{2mm}$=$\hspace{2mm}}{\textcolor{red!40!gray}{$\scriptscriptstyle#1$}}}{%
  \textcolor{red!40!gray}{$\scriptscriptstyle#2$}}}}
\newcommand{\egaprop}[1]{\underset{\textcolor{red!40!gray}{\scriptscriptstyle p.\ref*{prop_lim_#1}}}{\hspace{2mm}=\hspace{2mm}}}
\newcommand{\imphb}[2]{%
  \mathrel{\stackunder[2pt]{\stackon[4pt]{\hspace{2mm}$\Leftrightarrow$\hspace{2mm}}{\textcolor{red!40!gray}{$\scriptstyle#1$}}}{%
  \textcolor{red!40!gray}{$\scriptstyle#2$}}}}
\newcommand{\sol}[1]{\textcolor{solution}{{#1}}}
\newcommand{\solbf}[1]{\textcolor{solution}{\textbf{#1}}}
\newcommand{\solmath}[1]{\textcolor{solution}{\bm{#1}}}
\newcommand{\eme}{\textsuperscript{ème} } 
\newcommand{\er}{\textsuperscript{er} }
\newcommand{\ere}{\textsuperscript{ère} }  
%\newcommand{\pgdc}{\textrm{pgdc}}
\newcommand{\transp}[1]{{}^t \! #1} % pour noter la transposée d'une matrice
\newcommand{\CQFD}{\hfill\raisebox{-2mm}{$\blacksquare$}}
\newcommand{\dmo}[1]{\raisebox{#1mm}[0cm][0cm]{$\bm{-}$}} % (=DIvisionMOins)Pour les "moins" dans les tableaux des divisions euclidiennes : #1=ajustement de la hauteur (mm) | #2=ajustement du retrait (mm)
\NewDocumentCommand{\point}{O{-4} O{-4} O{-4} m}{\vspace{#3mm}\raisebox{#1mm}{\hspace{#2mm}\scriptsize\textcolor{pointscolor}{\textbf{/#4}}}} % 1: décalage vertical ; 2: décalage horizontal ; 3: espace interligne ; 4: nombre de points **** /!\ Il faut charger la package \usepackage{xparse} 
%"\NewDocumentCommand" permet de définir des commandes avec plusieurs arguments et leur valeur par défaut
%\NewDocumentCommand{\foocmd}{ O{default1} O{default2} m }{#1~#2~#3}
%                            %     ⤷ #1        ⤷ #2    ⤷ #3




%------------------
%---- Couleurs ----
%------------------
%\definecolor{cimpo}{gray}{0.7} % Couleur de fond pour des cadres
%\definecolor{crem}{gray}{0.4}  % Couleur pour du texte en remarque oral
\definecolor{cprop}{gray}{0.5}
%\definecolor{colhead}{RGB}{0, 112, 192} % Couleur Bleu clair
%\definecolor{colimportant}{RGB}{247 , 189 , 164}%{210, 76, 76}
%\newcommand{\colhead}{\color{Blue}}
%\definecolor{colhead}{RGB}{202, 25, 25}
\renewcommand*{\CancelColor}{\color{red}}
%\definecolor{cimpo}{gray}{0.7} % Couleur de fond pour des cadres
%\definecolor{crem}{gray}{0.4}  % Couleur pour du texte en remarque oral
\definecolor{cgrille}{gray}{0.3}
\definecolor{bleugris}{RGB}{76,100,125}
\definecolor{solution}{RGB}{213,35,35}
\definecolor{foncred}{RGB}{213,35,35}
\definecolor{foncblue}{RGB}{65,113,162}
\definecolor{foncgreen}{RGB}{100,162,65}
\definecolor{foncpurple}{RGB}{121,78,200}
\definecolor{foncyellow}{RGB}{187,182,56}
\definecolor{pointscolor}{RGB}{213,35,35}%{176,65,212}



%-------------------------------
%---- Raccourcis : tableaux ----
%-------------------------------
\newcolumntype{R}[1]{>{\raggedright}m{#1}}  % alignement à GAUCHE et spécification de la longueur de la colonne : R{4cm}
\newcolumntype{C}[1]{>{\centering\arraybackslash}m{#1}}
\newcolumntype{M}[1]{>{\centering\arraybackslash $}m{#1}<{$}}
\newcolumntype{S}{>{\small\it}l}
\newlength{\hauteurligne}
\newcommand{\hl}[1]{\setlength{\hauteurligne}{#1}\rule[(\hauteurligne*\real{-0.5})+0.1cm]{0cm}{\hauteurligne}}



%-------------------------------------------------------------------------------
%---- Raccourcis : Exemple(s) - retrait à droite et text italique sousligné ----
%-------------------------------------------------------------------------------
%---->>> Voir également les environnements
\newcommand{\exems}[1]{\hspace{5mm}\textsl{{Exemples:}} \hspace{#1mm}}
\newcommand{\exem}[1]{\hspace{5mm}\textsl{{Exemple:}} \hspace{#1mm}}



%------------------------------
%---- Raccourcis : Polices ----
%------------------------------
\newcommand{\comp}[1]{{\fontfamily{lmtt}\selectfont #1}}
\newcommand{\compbf}[1]{{\fontfamily{lmtt}\selectfont\bfseries #1}}
\newcommand{\compgui}[1]{{\fontfamily{lmtt}\selectfont\bfseries ``#1''}}



%---------------------------------------------------------
%---- Raccourcis : symboles mathématiques pour PYTHON ----
%---------------------------------------------------------
\newcommand{\pct}{\hspace{2pt}\%\hspace{2pt}}
\newcommand{\gui}{\textrm{"}}
\newcommand{\guic}{\comp{"}}
\newcommand{\mul}{\textrm{\hspace{1pt}*\hspace{1pt}}} % anciennement fois !
\newcommand{\puis}{\textrm{\hspace{1pt}**\hspace{1pt}}}
\newcommand{\id}{\hspace*{8mm}}		% Pour faire des identations lors de l'écriture de programmes
\newcommand{\idd}{\hspace*{16mm}}	% Pour faire des identations lors de l'écriture de programmes
\newcommand{\iddd}{\hspace*{24mm}}	% Pour faire des identations lors de l'écriture de programmes


%-----------------------------------------
%---- Raccourcis : Ecriture fonctions ----
%-----------------------------------------
\newcommand{\fonc}[5]{\renewcommand{\arraystretch}{1}$\begin{array}[t]{r r c l}
#1 : & \mathbb{#2} & \rightarrow & \mathbb{#3}\\
    & #4 & \mapsto & #5
\end{array}$}
\newcommand{\foncmi}[3]{\renewcommand{\arraystretch}{1}$\begin{array}[t]{r r c l}
#1 : & #2 & \mapsto & #3
\end{array}$}
%\newcommand{\fonction}[5]{\renewcommand{\arraystretch}{1}$\begin{array}[t]{r r c l}
%#1 : & #2 & \rightarrow & #3\\
%    & #4 & \mapsto & #5
%\end{array}$}
\newcommand{\fonction}[6]{\renewcommand{\arraystretch}{#6}$\begin{array}[t]{r r c l}
#1 : & #2 & \rightarrow & #3\\
    & #4 & \mapsto & #5
\end{array}$}
\newcommand{\function}[6]{\renewcommand{\arraystretch}{#6}\begin{array}[t]{r @{\hspace{1mm}}r c l}
#1 : & #2 & \rightarrow & #3\\
    & #4 & \mapsto & #5
\end{array}}



%-----------------------------
%---- Raccourcis : Limite ----
%-----------------------------
\newcommand{\lx}[2][x]{\lim\limits_{#1\rightarrow{#2}}}
\newcommand{\ly}[1]{\lim\limits_{y\rightarrow{#1}}}
\newcommand{\lh}[1]{\lim\limits_{h\rightarrow{#1}}}
\newcommand{\lt}[1]{\lim\limits_{t\rightarrow{#1}}}
\newcommand{\lin}[1]{\lim\limits_{n\rightarrow{#1}}}
\newcommand{\lk}[1]{\lim\limits_{k\rightarrow{#1}}}
\newcommand{\limx}[2]{\lim\limits_{x\rightarrow{#1}}\left(\displaystyle #2\right) }
\newcommand{\limy}[2]{\lim\limits_{y\rightarrow{#1}}\left(\displaystyle #2\right) }
\newcommand{\limh}[2]{\lim\limits_{h\rightarrow{#1}}\left(\displaystyle #2\right) }
\newcommand{\limt}[2]{\lim\limits_{t\rightarrow{#1}}\left(\displaystyle #2\right) }
\newcommand{\limn}[2]{\lim\limits_{n\rightarrow{#1}}\left(\displaystyle #2\right) }



%-----------------------------
%---- Raccourcis : Dérivées , Intégrales ----
%-----------------------------
\newcommand{\der}[1]{\displaystyle \left(#1\right)'}
\newcommand{\dx}[1][x]{\textrm{\textit{\hspace{2pt}d#1}}}



%--------------------------------
%---- Raccourcis : Géométrie ----
%--------------------------------
\renewcommand{\norme}[1]{\Vert#1\Vert}
\newcommand{\psca}{{\scriptstyle\hspace{0.6mm}\bullet\hspace{0.6mm}}} %produit scalaire
\newcommand{\pscar}{{\scriptstyle\hspace{0.3mm}\bullet\hspace{0.3mm}}} %produit scalaire espace réduit
\newcommand{\unite}[1]{\hspace{3pt}\textrm{#1}}
\newcommand\vv[1]{$\overrightarrow{#1}$}
\newcommand\vvert[2]{\begin{pmatrix*}[r]
#1 \\
#2
\end{pmatrix*}}
\newcommand\vvertnom[3]{\vec{#1}=\begin{pmatrix*}[r]
#2 \\
#3
\end{pmatrix*}}
\newcommand\vvvert[3]{\begin{pmatrix*}[r]
#1 \\
#2 \\
#3
\end{pmatrix*}}
\newcommand\vvvertnom[4]{\vec{#1}=\begin{pmatrix*}[r]
#2 \\
#3 \\
#4
\end{pmatrix*}}
\newcommand\vm[1]{\overrightarrow{#1}}
\newcommand{\pt}[2]{(#1\; ; #2)}
\newcommand{\ptt}[3]{(#1\; ; #2\; ; #3)}
\newcommand{\pttt}[4]{(#1\; ; #2\; ; #3\; ; #4)}
\newcommand{\ptnom}[3]{#1\hspace{0.5pt}(#2\; ; #3)}
\newcommand{\pttnom}[4]{#1\hspace{1pt}(#2\; ; #3\; ; #4)}



%-----------------------------------------
%---- Raccourcis : Tableau des Signes ----
%-----------------------------------------
\newcommand{\croi}{\boldsymbol{\nearrow}}
\newcommand{\decroi}{\boldsymbol{\searrow}}
\newcommand{\maxi}{\textrm{\footnotesize\bf max}}
\newcommand{\mini}{\textrm{\footnotesize\bf min}}
\newcommand{\palier}{\textrm{\footnotesize\bf palier}}
\newcommand{\inflex}{\textrm{\footnotesize\bf inflex.}}
\newcommand{\av}{\textrm{\footnotesize\bf A.V.}}
\newcommand{\nd}{\textrm{\footnotesize\bf n.d.}}
\newcommand{\conv}{\boldsymbol{\smile}}
\newcommand{\conc}{\boldsymbol{\frown}}



%-------------------------------------------
%---- Raccourcis : Systèmes d'équations ----
%-------------------------------------------
\newenvironment{systeme2}
	{
	\newcommand{\esp}{\hspace{2pt}}
	\newcommand{\p}{&+&}
	\renewcommand{\m}{&-&}
	\newcommand{\e}{&=&}
         $%
         \renewcommand{\arraystretch}{1.2}
            \left\lbrace\begin{array}{@{\esp}c@{\esp}@{\esp}c@{\esp}@{\esp}c@{\esp}@{\esp}c@{\esp}@{\esp}l}%
    }
	{
            \end{array}\right.%
         $%
	}
\newenvironment{systeme3}
	{
	\newcommand{\esp}{\hspace{2pt}}
	\newcommand{\p}{\esp &+&}
	\renewcommand{\m}{&-&}
	\newcommand{\e}{&=&}
         $
         \renewcommand{\arraystretch}{1.2}
            \left\lbrace\begin{array}{@{\esp}c@{\esp}@{\esp}c@{\esp}@{\esp}c@{\esp}@{\esp}c@{\esp}@{\esp}c@{\esp}@{\esp}c@{\esp}@{\esp}l}%
    }
	{
            \end{array}\right.%
         $%
	}
	


%-------------------------------
%---- Raccourcis : Matrices ----
%-------------------------------
\newenvironment{reducmatrice}[1]
	{
		\left(\begin{array}{@{}#1@{}}
	}
	{
		\end{array}\right)
	}
\newcommand{\ro}[1]{%
  \xrightarrow{\mathmakebox[\rowidth]{#1}}%
}
\newlength{\rowidth}% row operation width
\AtBeginDocument{\setlength{\rowidth}{3em}}

\newenvironment{elimination}
	{
		\everymath{\scriptstyle}
		\begin{array}{@{\scriptscriptstyle}c}
	}
	{
		\end{array}
	}


%-------------------------------------------------------------------------------
%---   FONCTION POUR FAIRE UNE CASCADE DE DIVISIONS POUR CHANGEMENT DE BASE  ---
%-------------------------------------------------------------------------------
\newcount\total
\newcount\lasttotal
\newcount\targetbase

\def\basetenconversiontable#1#2{%
    \begin{tikzpicture}[every node/.style={minimum width=1cm, minimum height=0.5cm}, x=1cm,y=0.5cm]
    %
    \total=#1%
    \targetbase=#2
    \def\newnumber{}
    %
    \pgfmathloop
    \ifnum\total<1
    \else
        %
        \ifnum\pgfmathcounter>1
            \node at (\pgfmathcounter, -\pgfmathcounter+1) (tmp) {\the\targetbase};
            \draw (tmp.north west) |- (tmp.south east);
            %
            \node at (\pgfmathcounter-1, -\pgfmathcounter) (tmp) {\pgfmathparse{int(\total*\targetbase)}\pgfmathresult};
            \draw (tmp.south west) -- (tmp.south east);
            %
            \pgfmathparse{int(\lasttotal-\total*\targetbase)}%
            \let\digit=\pgfmathresult
            \node at (\pgfmathcounter-1, -\pgfmathcounter-1) [text=red] {\digit};
            \edef\newnumber{\digit\newnumber}
        \fi
        %
        \ifnum\total<\targetbase
                \edef\newnumber{\the\total\newnumber}
            \node at (\pgfmathcounter, -\pgfmathcounter) [text=red]  {\the\total};
        \else
            \node at (\pgfmathcounter, -\pgfmathcounter) {\the\total};
        \fi
        \lasttotal=\total
        \divide\total by\targetbase
    \repeatpgfmathloop    
    \draw [->] (\pgfmathcounter-1,-\pgfmathcounter-1) -- ++(-0.5,0); 
    \node [anchor=west] at (1, -\pgfmathcounter-2) {$\overline{#1}_{10}=\overline{\newnumber}_{\the\targetbase}$};
    \end{tikzpicture}   
}


	

%--------------------------
%---- Labels Exercices ----
%--------------------------
%\newcommand{\exo}{\textbf{\textsl{Exercice \thechapter.\arabic{compteurexo}}} \vspace*{1mm}\\ \addtocounter{compteurexo}{1}}
\newcommand{\exo}{\refstepcounter{compteurexo}\textbf{\textsl{Exercice \thechapter.\arabic{compteurexo}}} \vspace*{1mm}\\}
\newcommand{\exonomme}[1]{\refstepcounter{compteurexo}\textbf{\textsl{Exercice \thechapter.\arabic{compteurexo} {\textsl{\footnotesize (#1)}}}}\vspace*{1mm}\\}
\newcommand{\exod}{\refstepcounter{compteurexo}\textbf{\textsl{Exercice \thechapter.\arabic{compteurexo}*}} \vspace*{1mm}\\}
\newcommand{\exodnomme}[1]{\refstepcounter{compteurexo}\textbf{\textsl{Exercice \thechapter.\arabic{compteurexo}* {\textsl{\footnotesize (#1)}}}} \vspace*{1mm}\\}
\newcounter{compteurexo}
%\setcounter{compteurexo}{1}
	
	

%-------------------------------------------
%---- Mises en formes pour les réponses ----
%-------------------------------------------
%%%%%%%%%%%%%%  CASE GRISE POUR REPONSES (avec espace après)
\newcommand{\caseesp}{\colorbox{gris15}{\makebox[4mm]{\mbox{\parbox[b][4mm]{\textwidth}{}}}}\;\;}
%%%%%%%%%%%%%%  CASE GRISE POUR REPONSES (sans espace après)
\newcommand{\case}{\colorbox{gris15}{\makebox[4mm]{\mbox{\parbox[b][4mm]{\textwidth}{}}}}}
%%%%%%%%%%%%%%  TRAITILLES POUR REPONSES AVEC PARAMETRE DE LONGUEUR
\newcommand{\rep}[1]{\raisebox{-2mm}{\makebox[#1cm]{\dotfill}}}
